\section{A common measure and a uniform modulus bound}
\label{cmp:section}

We now prove that the nonnegative ground function $\Phi$ has a fixed positive
constant-orbital occupation. The input is the conditional product law of the
unexposed trajectories from Section~\ref{group:section}, together with the good
lattice paths of Proposition~\ref{lat:paths}. We change selected trajectories
along one such path. Two assignments of their midpoints give the same observed
bath after different tagged particles are removed. The main estimate shows
that the second moment of this change of measure costs only encounters between
two lattice paths, rather than their lengths.

All constants in this section are independent of $D$, $K$, and the ground
vector once the earlier fixed parameters have been chosen. We first express
the occupation as a sum of affinities between two tagged position measures
and construct a common measure for them. We then prove the conditional
normalizer estimate that bounds this measure's second moments.

\subsection{Affinity and a common finite measure}

For finite nonnegative measures $P,Q$ on the same measurable space, define
their affinity by
\[
 \Aff(P,Q)=\int\sqrt{pq}\dd\xi,
 \qquad p=\frac{\dd P}{\dd\xi},\quad q=\frac{\dd Q}{\dd\xi},
 \quad \xi=P+Q.
\]
Using another dominating measure gives the same value. We will use three
elementary properties. If $P\ge P'$ and $Q\ge Q'$, then
$\Aff(P,Q)\ge\Aff(P',Q')$. For finite lists of measure pairs,
\begin{equation}\label{cmp:affinity-superadditive}
 \Aff\left(\sum_jP_j,\sum_jQ_j\right)
 \ge\sum_j\Aff(P_j,Q_j).
\end{equation}
Indeed, pointwise Cauchy--Schwarz gives
$\sum_j\sqrt{p_jq_j}\le\sqrt{(\sum_jp_j)(\sum_jq_j)}$.
Finally, applying the same probability kernel to both measures cannot decrease
affinity. To see this, form the joint measures with that kernel and then
condition a common dominating measure on the output. Conditional
Cauchy--Schwarz gives
$\E(\sqrt{pq}\mid\text{output})\le
\sqrt{\E(p\mid\text{output})\E(q\mid\text{output})}$.
Integration proves the assertion. Projection is a special case.

\begin{lemma}[A common finite measure]\label{cmp:holder}
Let $P,Q$ be finite nonnegative measures and let $\nu$ be a finite measure
of mass $m>0$, absolutely continuous with respect to both. If
\[
 \int\left(\frac{\dd\nu}{\dd P}\right)^2\dd P\le C_P,
 \qquad
 \int\left(\frac{\dd\nu}{\dd Q}\right)^2\dd Q\le C_Q,
\]
then $\Aff(P,Q)\ge m^2/\sqrt{C_PC_Q}$.
\end{lemma}
\begin{proof}
With densities $p,q,s$ relative to a common dominating measure, H\"older's
inequality with exponents $4,4,2$ gives
\[
 m=\int
 \left(\frac{s^2}{p}\right)^{1/4}
 \left(\frac{s^2}{q}\right)^{1/4}(pq)^{1/4}
 \le C_P^{1/4}C_Q^{1/4}\Aff(P,Q)^{1/2}.
\]
The ratios can be set to zero where their denominators vanish, since $s$ is
zero there. Squaring and rearranging proves the claim.
\end{proof}

\subsection{Cell measures and disjoint groups of tags}

Recall $V=K^3$, $D=N/V$, and the unit cells $B_v$. From this point onward,
$Y=(x_2,\ldots,x_N)$ again denotes the ordered bath, rather than a full
outer endpoint configuration. For ordered cells $v,w$, define
finite measures on $(x,y,Y)\in B_v\times B_w\times\Lambda_K^{N-1}$ by
\begin{equation}\label{cmp:cell-measures}
 \dd P_{vw}=V\Phi(x,Y)^2\dd x\dd y\dd Y,
 \qquad
 \dd Q_{vw}=V\Phi(y,Y)^2\dd x\dd y\dd Y.
\end{equation}
The occupation formula gives the exact identity
\begin{equation}\label{cmp:cell-affinity}
 B(\Phi)=\frac1{V^2}\sum_{v,w}\Aff(P_{vw},Q_{vw}).
\end{equation}

Fix a group index
$1\le g\le G_D:=\lfloor\ell D/(4M)\rfloor$ from
Section~\ref{group:section}. Its exposed data are $\cB$, its unexposed label
set is $\cI$, and its conditional path law is $\lambda_{\cI}$ from
Proposition~\ref{group:conditional-law}. Define measures
$\widehat P_g,\widehat Q_g$ that retain $\cB$ as well as $(x,y,Y)$.
For $\widehat P_g$, sample the data and the conditional unexposed paths, and
\emph{sum} the following contributions over all $i\in\cI$:
take $x=\omega_i(t)$, restrict to $x\in B_v$, order the remaining $N-1$
midpoints by a fresh uniform permutation to form $Y$, and sample $y$
uniformly in $B_w$. For $\widehat Q_g$, take the tagged midpoint as
$y\in B_w$ and sample the extra point $x$ uniformly in $B_v$.
The use of a sum, rather than a uniformly chosen tag, fixes their
normalization.

\begin{lemma}[Tag normalization]\label{cmp:tag-normalization}
For each ordered pair of cells,
\begin{equation}\label{cmp:group-affinity}
 \Aff(P_{vw},Q_{vw})
 \ge\frac1D\sum_{g=1}^{G_D}\Aff(\widehat P_g,\widehat Q_g).
\end{equation}
\end{lemma}
\begin{proof}
Put all groups on the original stationary sample, with all allocation
randomizations and a fresh uniform permutation of the full label set.
After deleting a tag, the restricted permutation uniformly orders its bath.
Each group's conditional construction is a disintegration of this common
sample with respect to its own exposed data. Thus the differing data spaces
are projected away before the following comparison; no conditional symmetry
of the unexposed-path law is asserted.
For every sampled path configuration, group label sets are disjoint. Their
sum of tag contributions is bounded by the sum over all $N$ trajectory
labels. The unrestricted midpoint vector has the symmetric density $\Phi^2$.
For each fixed tag label, uniform ordering of the bath therefore gives the
same density $\Phi(x,Y)^2$. Summing all tags gives $N$ times that density.
Since the extra target cell has volume one and $V/N=1/D$, projecting away
the data yields
\[
 \frac1D\sum_g\operatorname{proj}\widehat P_g\le P_{vw},
 \qquad
 \frac1D\sum_g\operatorname{proj}\widehat Q_g\le Q_{vw}.
\]
Use monotonicity, \eqref{cmp:affinity-superadditive}, and the increase of
affinity under projection. Its homogeneity supplies the factor $1/D$.
\end{proof}

The remaining task is to find, for each retained group and each of many
ordered cell pairs, a common measure of mass at least $1/2$ whose squared
likelihood integrals against $\widehat P_g$ and $\widehat Q_g$ are bounded
by one constant $C_2$. Lemma~\ref{cmp:holder} would then give affinity at
least $1/(4C_2)$ per group. The factor $1/D$ in
Equation~\eqref{cmp:group-affinity} can be paid because the number of
disjoint groups is proportional to $D$.

\subsection{Two midpoint assignments with the same observed bath}

Take $R=30$ in Proposition~\ref{lat:paths}, and fix an admissible
ordered cell pair $v,w$: its lift displacement $(n,d_2,d_3)$ satisfies
$\lceil K/10\rceil\le n\le\lfloor K/8\rfloor$ and $|d_i|\le n$.
Choose the good-site
parameters of Proposition~\ref{group:supply} to meet the bad-set hypothesis
of Proposition~\ref{lat:paths}. Let $\mathcal E_{vw}$ be the data-measurable
event supplied by Proposition~\ref{lat:paths}, of probability at least $1/2$. On this event draw a simple nearest-neighbor
path
\[
 z_0=v,z_1,\ldots,z_d=w
\]
through good sites. For each $k<d$, independently choose a uniform accessible
label $i_k$ anchored at $z_k$. The choices inspect only $\cB$ and the lattice
path. In particular, they do not inspect any latent trajectory. The selected
labels are distinct because their anchors are distinct.

We construct two path configurations that share their observed $(x,y,Y)$.
First sample all unselected unexposed trajectories from their exact marginal
under $\lambda_{\cI}$. For every vertex $z_k$, prepare a factor as follows.
The first copy needs label $i_k$ through this cell if $k<d$, and the second
copy needs label $i_{k-1}$ through it if $k>0$. Intersect the corresponding
good-target sets $T_{i_k}(z_k)$ and $T_{i_{k-1}}(z_k)$ from
Section~\ref{group:section}, using only the one required set at an endpoint.
The intersection, denoted by $A_k$, has volume at least $0.98$.
Sample $U_k$ uniformly in $A_k$, and then independently sample the one or
two required bridge trials through $U_k$, each conditioned on its individual
success against the frozen paths. All $d+1$ vertex factors are independent
before the additional conditioning below. Figure~\ref{cmp:shift-figure}
shows the two midpoint assignments.

Use $U_d$ as the extra point in the first copy and $U_0$ as the extra point
in the second copy. Impose all remaining hard-avoidance checks in both
copies: selected paths must avoid the unselected unexposed paths and one another
within their respective copies. Condition the product of vertex factors on
these checks, separately for each fixed collection of unselected paths.
The sampled unselected-path marginal therefore remains its original marginal
under $\lambda_{\cI}$.

We verify that this conditional construction is defined. A vertex factor
carries at most two paths, each with a marginal of the form $\Pi_i^{A_k}$
in Proposition~\ref{group:comparison}. Its paths are confined within distance
ten of their anchors, and there are at most $M$ unexposed labels per anchor.
Consequently each factor has at most $C_{\rm col}M$ possible opponents in the two
copies, for a fixed geometric constant $C_{\rm col}$. All opponents have distinct labels within the copy in which a check
is imposed. The deterministic-opponent estimate, summed over those checks,
gives
\begin{equation}\label{cmp:compatibility}
 \chi_D=C_{\rm col}M\eta_D\longrightarrow0
 \qquad(D\longrightarrow\infty,\ M\text{ fixed}).
\end{equation}
For any fixed regular and confined values of the other paths with their
recorded endpoints, the factor passes every incident check with probability
at least $1-\chi_D$.
This statement averages its common midpoint over $A_k$; it makes no claim at
a prescribed midpoint. It uses the two path marginals and a union bound,
so no independence of the paths within the factor is required. Inserting
the $d+1$ factors successively now bounds the two-copy conditioning
probability below by $(1-\chi_D)^{d+1}>0$ for sufficiently large $D$.

Tag $i_0$ in the first copy and $i_{d-1}$ in the second. In each copy the
tag and the extra point then give $x=U_0$, $y=U_d$. The selected bath
midpoints in both copies are exactly $U_1,\ldots,U_{d-1}$; the unselected
midpoints also agree. Choose a common uniform ordering of the shared midpoint instances; their
physical labels in the two copies may differ. In each copy separately,
this has the uniform ordering law determined by its paths and its tag.
This defines a single finite measure $\widehat\nu$ on $(\cB,x,y,Y)$ after
averaging the data, routes and labels, with zero mass off $\mathcal E_{vw}$.
Each conditional sampling law on this event is a probability measure, so
\begin{equation}\label{cmp:common-mass}
 \widehat\nu(\text{all})=\PP(\mathcal E_{vw})\ge\frac12.
\end{equation}

\begin{figure}[t]
\centering
\begin{tikzpicture}[x=2.45cm,y=.85cm,>=Stealth,
  every node/.style={font=\small}]
 \foreach \x/\txt in {0/{U_0},1/{U_1},2/{U_2},3/{U_d}} {
   \node at (\x,1.05) {$\txt$};
 }
 \node[anchor=east] at (-.4,0) {first copy};
 \node[anchor=east] at (-.4,-1.2) {second copy};
 \node[draw,rounded corners,fill=blue!8] (a0) at (0,0) {$i_0$ (tag)};
 \node[draw,rounded corners] (a1) at (1,0) {$i_1$};
 \node[draw,rounded corners] (a2) at (2,0) {$i_2$};
 \node at (3,0) {extra point};
 \node at (0,-1.2) {extra point};
 \node[draw,rounded corners] (b1) at (1,-1.2) {$i_0$};
 \node[draw,rounded corners] (b2) at (2,-1.2) {$i_1$};
 \node[draw,rounded corners,fill=blue!8] (b3) at (3,-1.2) {$i_{d-1}$ (tag)};
 \draw[->,gray] (a0) -- (b1);
 \draw[->,gray] (a1) -- (b2);
 \node at (2.48,-.52) {$\cdots$};
\end{tikzpicture}
\caption{The two midpoint assignments. Each departure label moves one cell
along the path in the second copy; the drawing suppresses the intermediate
cells before $U_d$. Removing the highlighted tags leaves the same bath
midpoints. Each vertical column is one proposal factor, with a common
midpoint and one path in each copy when both are required.}
\label{cmp:shift-figure}
\end{figure}

\subsection{Cancellation of conditional normalizers}

The midpoint construction has two path copies in each internal vertex
factor. To estimate its likelihood, designate one copy as primary and regard
the other as auxiliary. At an internal vertex $z_k$, $1\le k<d$, the
primary path has label $i_k$ in the first-copy view, whereas the auxiliary
path has label
$i_{k-1}$. Thus deleting a factor is not the same operation as deleting all
instances of one particle label. The following lemma keeps this distinction
explicit and isolates the normalizers that cancel when two routes are far
apart. Its coordinates and factors are abstract; we identify them with the
midpoint construction immediately after the proof.

The next lemma concerns finitely many coordinates in standard Borel spaces,
each with a probability reference $P_i$. All auxiliary spaces below are
standard Borel as well. A \emph{check} will mean an indicator imposing a constraint
on one or two coordinates. The joint reference law $\lambda$ is the product
$\bigotimes_iP_i$ conditioned on a finite collection of checks. An additional
independent coordinate is allowed; it simply has no checks.

A change indexed by a set $S$ leaves the marginal outside $S$ unchanged.
Inside $S$, it starts from independent factors $\Pi_{S,i}$, $i\in S$.
Each factor includes a \emph{primary coordinate} in the space of $P_i$ and
possibly auxiliary coordinates. Its primary marginal is $g_{S,i}P_i$, with
\begin{equation}\label{cmp:factor-density}
 0\le g_{S,i}\le G,
 \qquad \int g_{S,i}\dd P_i=1.
\end{equation}
The factors are fixed before the outside coordinates are sampled. The product
is conditioned on specified checks, including every reference check involving
the primary coordinates in $S$. Checks involving only fixed outside
coordinates are already satisfied almost surely under the unchanged outside
marginal; otherwise the conditional proposal would not be defined there.
Omit these redundant checks. Every remaining proposal check therefore
involves a coordinate carried by a changed factor. Auxiliary coordinates
are then forgotten.
The resulting primary likelihood relative to $\lambda$ is denoted by $L_S$.

For each primary and auxiliary coordinate, fix a measurable allowed set
on which all its relevant reference and proposal marginals are
concentrated. These sets, like the proposal factors, are fixed before
sampling outside coordinates. A configuration is \emph{individually
permitted} when each coordinate lies in its allowed set; this imposes no
pairwise checks. The compatibility hypothesis below is uniform on these
sets, including configurations that fail some pairwise checks.

We make explicit the deletion convention for this construction. When some
factors are removed, remove their coordinates and all checks incident to them.
Do not replace a removed auxiliary coordinate by an original reference
coordinate. Coordinates fixed outside the sets under discussion remain as
obstacles in all the checks in which they occur. Coordinates here are path
instances in a specified copy. If an auxiliary with label $j$ belongs to a
retained factor, it remains integrated in that factor even when a different
factor carrying the primary label $j$ is deleted. It is never replaced by
the original path with label $j$.

\begin{lemma}[Cancellation away from an overlap]\label{cmp:cancellation}
Consider two changes $S,T$ of the preceding kind, with $G\ge1$, and put
$U=S\cup T$. Suppose that the following conditions hold.
\begin{enumerate}[label=\textup{(\roman*)}]
 \item There is $\chi\in[0,1)$ such that inserting any coordinate of $U$
 from its reference $P_i$, or any factor of either proposal from its
 $\Pi_{S,i}$ or $\Pi_{T,i}$, satisfies all checks incident to that insertion
 with probability at least $1-\chi$. This bound is uniform in every fixed
 choice of the other individually permitted coordinates. It remains valid
 when any checks or coordinates have been deleted.
 \item There are decompositions $S=C_S\sqcup F_S$ and
 $T=C_T\sqcup F_T$, with $F_S\cap T=F_T\cap S=\varnothing$.
 No coordinate carried by an $F_S$ factor has a check with any coordinate
 carried by a $T$ factor, or with the original primary coordinates indexed
 by $T$; the corresponding assertion holds with $S,T$ interchanged.
 The reference coordinates in $F_S,F_T$ also have no mutual checks.
\end{enumerate}
Then
\begin{equation}\label{cmp:cancellation-bound}
 \int L_SL_T\dd\lambda
 \le G^{|S\cap T|}(1-\chi)^{-2(|C_S|+|C_T|)}.
\end{equation}
It is enough that the assumptions hold for almost every fixed admissible
configuration outside $U$.
\end{lemma}

\begin{proof}
Put $a_\chi=1-\chi$ and $c=|C_S|+|C_T|$. Fix the primary coordinates
outside $U$. All normalizers below are conditional on these fixed values;
we first work on a configuration for which this conditional law is defined.

For any selected set, its reference normalizer is the integral of its
reference-check indicator against its product references. In particular,
write $Z_U$ for this integral over $U$. When the coordinates outside $S$
are fixed, write $Z_S$ for the corresponding integral over $S$. Thus $Z_S$
can depend on coordinates in $T\setminus S$. Write $Y_S$ for the integral
of the proposal-check indicator $J_S$ against $\bigotimes_{i\in S}\Pi_{S,i}$,
with the same outside values fixed.

Disintegrate each factor in the form
\[
 \Pi_{S,i}(\dd x_i\,\dd\zeta_i)
   =P_i(\dd x_i)K_{S,i}(x_i,\dd\zeta_i),
 \qquad K_{S,i}(x_i,\text{all})=g_{S,i}(x_i).
\]
Here $\zeta_i$ denotes all its auxiliary coordinates. Define
\[
 H_S(x_S;x_{S^c})
  =\int J_S\prod_{i\in S}K_{S,i}(x_i,\dd\zeta_i).
\]
On the reference-admissible configurations, the exact likelihood is
\begin{equation}\label{cmp:exact-likelihood}
 L_S=\frac{Z_S}{Y_S}H_S.
\end{equation}
Indeed, the conditional reference density is its check indicator divided by
$Z_S$, whereas the new primary density is $H_S/Y_S$, both relative to
$\bigotimes_{i\in S}P_i$. The outside marginal is identical. The inclusion
of primary reference checks ensures that the new density is zero on
reference-forbidden configurations.

Now retain only the factors indexed by $F_S$, together with coordinates fixed
outside $U$, using the stated deletion convention. Denote their reference
normalizer, proposal normalizer and primary proposal density by
$Z_{F_S},Y_{F_S},H_{F_S}$. Define the three corresponding $F_T$ quantities
using the factors from the $T$ change. These quantities depend only on the
outside of $U$ and, for $H_{F_S},H_{F_T}$, on their own primary coordinates.
They do not depend on the omitted near coordinates.

Successive insertion using assumption~\textup{(i)}, and noninteraction in
assumption~\textup{(ii)}, give
\begin{align}
 Z_U&\ge a_\chi^c Z_{F_S}Z_{F_T},
       \label{cmp:union-normalizer}\\
 Z_S&\le Z_{F_S},\qquad
 Y_S\ge a_\chi^{|C_S|}Y_{F_S},
       \label{cmp:single-normalizers}\\
 Z_T&\le Z_{F_T},\qquad
 Y_T\ge a_\chi^{|C_T|}Y_{F_T}.
       \label{cmp:other-normalizers}
\end{align}
For the first inequality, the reference integral over the two far sets
factorizes; add every remaining distinct coordinate with its uniform
compatibility bound. There are at most $c$ such coordinates. For the upper
bound on $Z_S$, delete the checks involving $C_S$ and integrate those reference
coordinates freely. Coordinates in $T\setminus S$ have no checks with the
far set. For its proposal lower bound, start with the reduced $F_S$ integral
and add the $C_S$ factors. The fixed original coordinates in
$T\setminus S$ still do not meet the far factors. Each addition retains at
least the factor $a_\chi$. Applying these two arguments to the $T$ change proves the
last line. These statements are pointwise in the opposing coordinates;
they do not assert that $Z_S$ or $Y_S$ is constant in them.

The insertion bounds also prove positivity of the divisors. In particular,
integrating coordinates one at a time gives a positive product lower bound
for every reduced normalizer. The full reference conditional is defined on
the outside configurations under consideration.

Deleting proposal checks involving a near factor gives
\[
 H_S\le H_{F_S}\prod_{i\in C_S}g_{S,i},
 \qquad
 H_T\le H_{F_T}\prod_{i\in C_T}g_{T,i}.
\]
After dropping the reference admissibility indicator in an upper bound,
the two far primary integrals equal $Y_{F_S}$ and $Y_{F_T}$. The remaining
coordinates carry only the displayed marginal densities. A coordinate
occurring once integrates to one. A coordinate occurring in both lists
costs at most $G$, by \eqref{cmp:factor-density}. Therefore
\begin{equation}\label{cmp:kernel-product}
 \int H_SH_T\ind_{\mathrm{admissible}}
              \prod_{i\in U}P_i(\dd x_i)
 \le Y_{F_S}Y_{F_T}G^{|S\cap T|}.
\end{equation}
The conditional reference law on $U$ has normalizer $Z_U$. Combine
\eqref{cmp:exact-likelihood}--\eqref{cmp:kernel-product}. The two pointwise
ratio bounds contribute
$a_\chi^{-c}Z_{F_S}Z_{F_T}/(Y_{F_S}Y_{F_T})$; division by $Z_U$
contributes at most $a_\chi^{-c}/(Z_{F_S}Z_{F_T})$. All four far
normalizers cancel, leaving $G^{|S\cap T|}a_\chi^{-2c}$.
Finally integrate the outside coordinates. This proves
\eqref{cmp:cancellation-bound}.
\end{proof}

\subsection{The likelihood cost is confined to encounters}

Fix the data, the route and the selected labels. Regard either copy as the
primary one. The reference for that copy is
$\lambda_{\cI}\otimes P_*$, where $P_*$ is the uniform law of the extra
point in its endpoint cell. Let
\[
 S=\{i_0,\ldots,i_{d-1}\}\cup\{*\}.
\]
The objects in Lemma~\ref{cmp:cancellation} now have the following meaning:
\begin{center}
\begin{tabular}{@{}p{.24\linewidth}p{.70\linewidth}@{}}
\toprule
Primary coordinate & The chosen copy's path, or its endpoint extra point.\\
Auxiliary coordinate & The other copy's path or extra point through the same midpoint.\\
Checks & Collision exclusions within each copy.\\
\bottomrule
\end{tabular}
\end{center}
Index a vertex factor by its primary label, using $*$ for the extra point;
its center is its midpoint cell $z_k$. The two paths in one factor may have
different labels, and each copy contributes at most one path. The allowed
set for each path consists of the regular, confined, frozen-avoiding paths
with its own recorded endpoints. Each extra point lies in its target cell.
These are individual restrictions; the remaining collision checks still
couple different factors.

By Proposition~\ref{group:comparison}, the primary marginal of each path
factor has density at most $G$ relative to $P_i$, and integrates to one.
For the extra point the density is $\ind_{A_k}/|A_k|$ relative to $P_*$,
which is at most $1/0.98$. Enlarge the fixed $G$ to be at least two, covering
this factor as well. These marginal densities and their auxiliary kernels
depend on the exposed data and the route and label choices; they do not
depend on the unselected unexposed paths.

The union bound proving Equation~\eqref{cmp:compatibility} verifies
hypothesis~\textup{(i)} for each proposal factor. The same argument with
$P_i$ in place of $\Pi_i^{A_k}$ verifies it for insertion of a selected
reference path, by Proposition~\ref{group:comparison}. The reference extra
point has no checks. These estimates are uniform over all individually
permitted opponent paths, whether or not their labels are accessible, and
deleting checks can only improve them. Thus the bound remains available
for every reduced normalizer in Lemma~\ref{cmp:cancellation}.

Take two independent route and label choices at the same data and endpoint
pair, and denote their selected index sets by $S,T$. For the chosen primary
side, declare a factor of $S$ near if its center is at lattice supremum
distance at most $R$ from some center of $T$, and declare every other factor
far. Define $C_S,F_S,C_T,F_T$ accordingly. A shared real label has centers
at distance at most two in the two assignments, so it is near; the shared
extra index has the same endpoint center. Consequently the far index sets
are disjoint from the opposing lists.

All paths carried by a factor lie within distance eleven of its center.
The original path of an opposing real label also lies within distance eleven
of that opposing factor's center. Since $R=30$, far factors cannot interact
with any coordinate carried by the opposite list or with those original
paths. This includes the auxiliary-copy paths. Far sets therefore meet the
separation assumptions of Lemma~\ref{cmp:cancellation}.

For clarity, a reduced far proposal retains exactly the paths carried by
its far vertex factors and the fixed paths outside $S\cup T$. It deletes
the omitted near factors in both copies. It does not put their original
trajectories back as unselected obstacles. With this convention all the
normalizers in \eqref{cmp:union-normalizer}--\eqref{cmp:kernel-product}
are the ones arising from the present construction.

Write $L_S,L_T$ for the two primary likelihoods. If
\[
 J_R=\#\{(j,k):d_\infty(z_j,z'_k)\le R\},
\]
then $|C_S|+|C_T|\le2J_R$. Lemma~\ref{cmp:cancellation} gives, at the
fixed data,
\begin{equation}\label{cmp:two-choice-bound}
 \int L_SL_T\dd(\lambda_{\cI}\otimes P_*)
 \le G^{|S\cap T|}(1-\chi_D)^{-4J_R}.
\end{equation}
In particular, no factor is paid for a long portion of one route that stays
away from the other.

\subsection{Averaging labels and forgetting the paths}

Condition on the data and the two lattice paths before choosing their labels.
Each common departure cell has at least $M/2$ accessible labels. The two
uniform choices at that cell match with probability at most $2/M$, and the
match events are independent between distinct common cells. Real labels
anchored at different cells cannot agree. The shared extra point contributes
one factor $G$, so
\begin{equation}\label{cmp:label-average}
 \E_{\mathrm{labels}}G^{|S\cap T|}
 \le G\exp\left(\frac{2(G-1)}M J_R\right).
\end{equation}
Choose the fixed integer $M$ sufficiently large for the good-site supply
and for
\begin{equation}\label{cmp:M-choice}
 \frac{2(G-1)}M\le\frac\kappa2.
\end{equation}
Then increase $D_0$ so that, for every $D\ge D_0$ in the allowed $\alpha$
window, all earlier estimates hold and
\begin{equation}\label{cmp:D-choice}
 -4\log(1-\chi_D)\le\frac\kappa2,
 \qquad
 \frac{G_D}{D}\ge\frac\ell{8M}.
\end{equation}
The first requirement is possible by \eqref{cmp:compatibility}; the second
follows once $\ell D/(4M)\ge2$. These choices do not change $G$, $\kappa$,
or $C_{\mathrm{path}}$.

The following use of an averaged likelihood and a two-route second moment
also occurs in the point-reassignment argument of
Peres and Sly~\cite[Proposition~2.1]{PeresSly2014}. Here
Lemma~\ref{cmp:cancellation} supplies the additional conditional-normalizer
estimate required by the hard-core path law.

Let $\overline L$ be the likelihood averaged over the conditional route and
label choices, with value zero off $\mathcal E_{vw}$. Tonelli's theorem
expresses its square integral using two independent such choices. Combining
\eqref{cmp:two-choice-bound}--\eqref{cmp:D-choice} and the annealed encounter
bound in Proposition~\ref{lat:paths} gives
\begin{equation}\label{cmp:latent-second-moment}
 \E_{\cB}\int\overline L^{2}
           \dd(\lambda_{\cI}\otimes P_*)
 \le G\E\bigl[\ind_{\mathcal E_{vw}}e^{\kappa J_R}\bigr]
 \le C_2,
 \qquad C_2:=G C_{\mathrm{path}}.
\end{equation}
Here the two route randomizations share the same data; the expectation over
the data is taken only once. The event is not conditioned to have probability
one. Its indicator is part of both likelihoods, and its square is itself.
The calculation applies separately with either copy primary, with the same
constant $C_2$.

We next check carefully that tagging and observing the bath do not increase
this bound. At fixed data, split $\overline L=\sum_{i\in\cI}f_i$, where
$f_i$ contains exactly the choices designating $i$ as the primary tag.
The functions are nonnegative. In the first copy $f_i$ is supported where
$\omega_i(t)\in B_v$; in the second it is supported where
$\omega_i(t)\in B_w$. Since
\[
 \sum_i f_i^2\le\left(\sum_i f_i\right)^2,
\]
the measure that retains the tag index has squared likelihood integral at
most $C_2$ relative to the \emph{sum} of the tag-restricted reference
measures on the disjoint copies indexed by $i$. Now apply the observation kernel that orders the bath and retains
$\cB,x,y,Y$, and forget the tag and all the paths. If a finite measure has
density $f$ relative to a reference $P$, the density after applying a common
kernel is the conditional average of $f$ in the corresponding joint
reference measure. Conditional Cauchy--Schwarz bounds its squared integral
by $\int f^2\dd P$. This remains valid for finite reference measures,
by normalization when their mass is nonzero.

For either primary copy, the resulting measure is the common observed
$\widehat\nu$ already constructed. Its corresponding reference is precisely
$\widehat P_g$ or $\widehat Q_g$. We have proved
\begin{equation}\label{cmp:observed-second-moments}
 \int\left(\frac{\dd\widehat\nu}{\dd\widehat P_g}\right)^2
                \dd\widehat P_g\le C_2,
 \qquad
 \int\left(\frac{\dd\widehat\nu}{\dd\widehat Q_g}\right)^2
                \dd\widehat Q_g\le C_2.
\end{equation}
Together with \eqref{cmp:common-mass}, Lemma~\ref{cmp:holder} now gives
\begin{equation}\label{cmp:one-group-gain}
 \Aff(\widehat P_g,\widehat Q_g)\ge\frac1{4C_2}.
\end{equation}
This lower bound is uniform in the group index and in every admissible
ordered pair of cells.

\begin{proposition}[Uniform modulus occupation]\label{cmp:modulus}
After the fixed choices above, there are constants $D_0<\infty$ and
\begin{equation}\label{cmp:c-star}
 c_*:=\frac{\upsilon_0\ell}{32M G C_{\mathrm{path}}}>0,
 \qquad \upsilon_0=\frac1{2000}
\end{equation}
such that, whenever $D\ge D_0$, $\alpha=Dr\in[\alpha_0/2,2\alpha_0]$,
and $K$ is sufficiently large, every normalized nonnegative bosonic ground
function satisfies $B(\Phi)\ge c_*$. The constant $c_*$ is unchanged if
the lower threshold on $D$ is subsequently increased. The volume thresholds
can be taken locally uniform for $D$ in compact subintervals of
$[D_0,\infty)$.
\end{proposition}
\begin{proof}
For large $K$, at least $\upsilon_0V^2$ ordered cell pairs meet the displacement
condition of Proposition~\ref{lat:paths}. For each such pair,
Lemma~\ref{cmp:tag-normalization}, \eqref{cmp:D-choice} and
\eqref{cmp:one-group-gain} give
\[
 \Aff(P_{vw},Q_{vw})
 \ge\frac{G_D}{4D C_2}
 \ge\frac\ell{32M C_2}.
\]
Sum these nonnegative contributions in \eqref{cmp:cell-affinity} to obtain
$B(\Phi)\ge\upsilon_0\ell/(32M C_2)$, which is \eqref{cmp:c-star}.
All ingredients are uniform in the ground vector. The threshold and
local-uniformity assertions follow from the preceding path and group
estimates and Proposition~\ref{lat:paths}. No parameter entering $c_*$ is
chosen after $D_0$.
\end{proof}

The modulus occupation has now been bounded below by a fixed number. To
treat the original complex ground vector, it remains to show that a small
part of each one-coordinate slice contains every possible phase obstruction.
That estimate may require a larger lower bound on $D$; Proposition~\ref{cmp:modulus}
permits that increase without changing the gain.
